\clearpage
\fancyhead[L]{\small CONNAISSANCE DES TEMS.}
\hypertarget{native-643-lower}{}
\begin{center}\small Göttingische gelehrte Anzeigen. Stück 80. Seiten 793..796. 17. Mai 1824.\\[3pt]\rule{.72\linewidth}{.3pt}\end{center}

\emph{Connaissance des tems ou des mouvemens célestes à l'usage des astronomes et des navigateurs pour l'an \textup{1826}. Publiée par le bureau des longitudes.} Paris 1823. Bei Bachelier. Die stehenden Artikel 216 Seiten, die Zusätze 110 Seiten. Großoktav.

Zie Zusätze zu diesem Jahrgang bestehen aus den folgenden Artikeln. Übersicht der astronomischen Beobachtungen, die \textsc{Gambart} d.\,J. in den Jahren 1820..1822 an der Sternwarte von Marseille angestellt hat. Die Sternwarte von Marseille war bisher mit wenigen und mittelmäßigen Instrumenten ausgestattet, und Herr \textsc{Gambart} hat sich daher mit Recht auf Beobachtungen beschränkt, die diesem Umstand angemessen sind. Der fast immer klare Himmel von Marseille lädt besonders zur Suche nach neuen Kometen ein. Das Glück, das früher den dadurch berühmt gewordenen \textsc{Pons} so sehr begünstigt hat, scheint auch \textsc{Gambart} treu zu bleiben. Von den drei Kometen des Jahres 1822 hat er einen als Erster und einen weiteren zumindest ohne Kenntnis der früheren Entdeckung in Marlia selbst entdeckt. \textsc{Gambart}s wertvolle Beobachtungen dieser beiden Kometen sind hier so ausführlich abgedruckt, wie man es sich nur wünschen kann. Auch am zweiten Kometen des Jahres 1822 hat \textsc{Gambart} zwei Beobachtungen gemacht, die umso wertvoller sind, als diesen Kometen außer ihm nur \textsc{Caturegli} in Bologna beobachtet hat. Die übrigen Beobachtungen betreffen Verfinsterungen von Jupiterstrabanten, Sternbedeckungen und die Mondfinsternis vom 2. August 1822. --- Über die Verteilung der Wärme in einem homogenen Ring, der überall gleich dick ist und sich in einem Raum befindet, in dem an verschiedenen Punkten unterschiedliche Temperaturen herrschen [von \textsc{Poisson}]. Wegen seiner feinen mathematischen Behandlung interessiert dieser schöne Aufsatz hauptsächlich den Geometer. Er ist aber auch für die praktische Astronomie von Interesse, da sich Kreisinstrumente bei Beobachtungen nicht selten unter Umständen befinden, wie sie hier zugrunde gelegt \SourceLeaf{644} sind, und es wünschenswert ist, eine genaue Vorstellung von der daraus entstehenden Verformung der Instrumente zu gewinnen. --- Über die Geschwindigkeit des Schalls [von \textsc{Poisson}]. Eine ebenso gründliche wie klare Darstellung der Theorie dieses Gegenstands, die zuerst von \textsc{Laplace} vollendet wurde. Verwendet man das Verhältnis der spezifischen Wärme der Luft bei konstantem Druck zu ihrer spezifischen Wärme bei konstantem Volumen, wie es sich aus dem interessanten Versuch von \textsc{Clement} und \textsc{Delorme} ergibt, so gelangt man zu einer beinahe vollkommenen Übereinstimmung zwischen der Theorie der Schallgeschwindigkeit und den empirischen Ergebnissen. --- Beobachtungen des ersten und dritten Kometen des Jahres 1822 an der Pariser Sternwarte. Den letzteren hatte Herr \textsc{Bouvard} selbst entdeckt, wenn auch einige Tage später als \textsc{Pons} und \textsc{Gambart}. --- Ergebnisse aus Beobachtungen mit unveränderlichen Pendeln, berechnet von Herrn \textsc{Mathieu}. Der Schiffsleutnant \textsc{Duperrey}, der in den Jahren 1822 und 1823 mit der Korvette Lacoquille eine Reise um die Welt machte, führte zwei unveränderliche Pendel mit sich. Ihre Schwingungen wurden vor der Abreise in Paris und von Herrn \textsc{Duperrey} bei einer Landung auf den Malouinen beobachtet. Da deren südliche Breite nur um $2^\circ\,41'$ größer ist als die nördliche Breite von Paris, lässt sich aus diesen Versuchen, die im Übrigen mit großer Sorgfalt erörtert wurden, eigentlich nur folgern, dass sie mit einer Abplattung von $\frac{1}{308}$ vereinbar sind. Wichtiger sind die hier ebenfalls berechneten Versuche mit einem unveränderlichen Pendel, die \textsc{Brisbane} und \textsc{Rümker} zuerst in London und danach in Paramatta auf Neuholland angestellt haben und aus denen hier eine Abplattung von $\frac{1}{206}$ berechnet wird. --- Über die Wirkung des Mondes auf die Atmosphäre, von Herrn \textsc{de Laplace}. Seit 1815 wird die Barometerhöhe an der Pariser Sternwarte täglich viermal beobachtet, nämlich vormittags um 9 Uhr, mittags, nachmittags um 3 Uhr und abends um 9 Uhr. Die tägliche Schwankung des Barometerstands hebt sich so deutlich von den zufälligen Schwankungen ab, dass die Differenz zwischen dem Stand vormittags um 9 Uhr und dem Stand nachmittags um 3 Uhr im Mittel jedes einzelnen Monats während sechsjähriger Beobachtungen kein einziges Mal ihr Vorzeichen geändert hat (im Mittel beträgt die Differenz 0.801 Millimeter, wobei der Vormittagsstand höher liegt). Das gab Anlass zu dem Versuch, aus einer geeigneten Auswahl dieser Beobachtungen auch den Einfluss des Mondes auf den Barometerstand zu ermitteln. Zu diesem Zweck wurden die Beobachtungen in der Nähe der Syzygien mit denen in der Nähe der Quadraturen verglichen. Der Vergleich ergab eine kleine periodische Ungleichheit mit einer Periode von einem halben Mondtag. Ihr Maximum beträgt nur 0.0272 Millimeter und tritt am Tag der Syzygien um $3^{\rm u}18^{\rm m}36^{\rm s}$ ein. Da sich in den numerischen Teil dieser Untersuchung einige Ungenauigkeiten eingeschlichen zu haben scheinen, sah sich der Rezensent dadurch und durch das Interesse des Gegenstands veranlasst, die Rechnungen teils genauer, teils nach anderen Methoden zu wiederholen. Das Ergebnis ist jedoch fast dasselbe wie das von \textsc{Laplace}. Der Rezensent findet für das Maximum 0.0269 Millimeter, für dessen Zeitpunkt $3^{\rm u}11^{\rm m}49^{\rm s}$ und für den bei der ersten Bestimmung zu befürchtenden mittleren Fehler 0.0264 Millimeter. --- Den Abschluss dieses Jahrgangs bildet ein kurzer Aufsatz von Herrn \textsc{Puissant} über die Berechnung der mit einem Repetitionstheodoliten gemessenen Azimute. Das hier vorgeschlagene Verfahren ist bereits bekannt und stimmt im Wesentlichen völlig mit dem Verfahren überein, das Herr \textsc{Soldner} in den Denkschriften der Münchner Akademie für 1813 abgehandelt hat und zu dem man auch unsere Blätter, Jahrgang 1815, S.\ 449 ff.\ [S.\ 587 d.\,B.], nachschlagen kann.
